If activity rungs cannot persist information, the weights must: they
survive silence by construction (they are only changed by plasticity
events), and the record shows they are the substrate's persistent store.
The plasticity-as-memory audit \cite{synmem} established this with four
independent verdicts.

\paragraph{Persistence (G1, G2).}
Cross-run training (A-only vs.\ C-only, same seed, identical initial
weights) produces distinct persistent configurations: cosine 0.119,
selectivity \(\pm 0.8\)--\(0.9\), 91\% leave-one-pair-out run
identification. Selectivity develops monotonically through S1 and
\emph{continues to develop in silence} (\(+0.903 \to +0.914\), driven by
pruning of residual non-A mass --- 617\(\to\)309 live afferents),
while the matrix itself changes only \(\sim\)1.4\% across 20 s of
silence. Structure persists; activity decays (activity rungs decay
3.5--7\(\times\) in the same window: \(u\)-norm 7.999\(\to\)2.282,
internal spikes 479/s \(\to\) 71/s).

\paragraph{Create/maintain mechanics (G4).}
The causal records (E15--E17, SDE series) attribute organization to a
specific subset of mechanisms: M2 normalization (per-neuron \(t_e\)
budget) is \emph{required} for the separated regime (M2-off destroys
separation; E16), while STDP is not individually necessary for the
afferent-selection trajectory (STDP-off reproduces it near-verbatim,
E16), and M3/M4 topology dynamics are not necessary either (E17).
Maintenance is M2 cadence plus M4 pruning; passive decay
(\(10^{-6}\)/tick) is the dominant single eraser of first-block mass
(SDE-C2: decay-off preserves the first cohort in 6/6 blocked runs;
residual skew \(\approx 0.07\)--\(0.09\) remains). Retrieval is a
dynamics function: the REV1 phase-reversal experiment re-expresses an
early-established configuration with zero plasticity events and
\(\le 3\%\) weight change (E12--E17) --- forward dynamics through
stored weights.

\paragraph{What is not established at this layer (G8).}
Re-expression of a stored configuration by a later input exists at the
association/trace level; \emph{episode-specific} reinstatement under
alternation does not --- the alternating store superposes
(§\ref{sec:interference}) and no later \(A\)-presentation reinstates an
A-configuration distinguishable from the adjacent C-configuration.

\begin{quote}
\textbf{Demonstrated:} synaptic structure is the persistent store;
creation requires budget normalization; retention is bounded by decay
and order; retrieval-by-forward-dynamics exists for trace-level
configurations.\\[1pt]
\textbf{Supported:} the information bottleneck is shared by every
persistent substrate tested (weights and \(u\) show the same
superposition signature), hence the limiting factor is the shared
additive store under a shared budget, not the substrate identity
\cite{synmem}.
\end{quote}